\begin{resumo}
Organizações de desenvolvimento de software geram grandes volumes de dados de processo por meio de sistemas de gerenciamento de \textit{issues}, repositórios de controle de versão, \textit{pipelines} de integração contínua e plataformas de revisão de código. Apesar dessa extensa instrumentação, a previsão do comportamento de processos de desenvolvimento de software ainda se baseia predominantemente em métricas agregadas e métodos heurísticos de estimativa que tratam os fluxos de trabalho de desenvolvimento como caixas-pretas. Como consequência, as abordagens existentes fornecem uma compreensão limitada da dinâmica estrutural e da variabilidade dependente do caminho que caracterizam os processos reais de desenvolvimento de software. Esta tese investiga se modelos derivados de mineração de processos podem fornecer uma base mais informativa para a previsão probabilística de processos de desenvolvimento de software. Em particular, a pesquisa testa a hipótese de que características derivadas do processo e modelos estocásticos explícitos podem melhorar o desempenho preditivo em comparação com abordagens baseadas exclusivamente em métricas tradicionais de repositórios de software. Para abordar esse problema, a tese propõe o PATHCAST (PATH-aware foreCASTing), um método computacional que integra mineração de processos, modelagem por cadeias de Markov, simulação de Monte Carlo e aprendizado de máquina em um \textit{pipeline} analítico unificado. A abordagem transforma dados heterogêneos provenientes de repositórios de software em logs de eventos estruturados, a partir dos quais são descobertos modelos de processo que representam os fluxos reais de desenvolvimento. Esses modelos são então utilizados para derivar modelos estocásticos de transição de estados que permitem a simulação probabilística de trajetórias de execução do processo e de suas distribuições de resultados. A pesquisa segue a metodologia \textit{Design Science Research}, na qual o método proposto é implementado e avaliado empiricamente com base em dados históricos de projetos de software. A avaliação compara o desempenho preditivo do PATHCAST com modelos de referência baseados em métricas de repositórios e em técnicas convencionais de previsão. Espera-se que os resultados demonstrem que a incorporação explícita da estrutura do processo em modelos de previsão melhore tanto a precisão preditiva quanto a interpretabilidade, contribuindo, assim, para o avanço da inteligência de processos aplicada à engenharia de software.
	\palavraschave{Mineração de Processos, Previsão Probabilística, Cadeias de Markov, Simulação de Monte Carlo, Aprendizado de Máquina, Design Science Research}
\end{resumo}
